\documentclass{article}

\usepackage{amsmath,amssymb}
\usepackage{xurl}
\usepackage[hidelinks]{hyperref}
\setlength{\emergencystretch}{2em}

% Shared commands for notes in docs/paper-gaps/.

\DeclareUnicodeCharacter{2080}{\ifmmode{}_{0}\else\textsubscript{0}\fi}
\DeclareUnicodeCharacter{2081}{\ifmmode{}_{1}\else\textsubscript{1}\fi}
\DeclareUnicodeCharacter{2082}{\ifmmode{}_{2}\else\textsubscript{2}\fi}
\DeclareUnicodeCharacter{2099}{\ifmmode{}_{n}\else\textsubscript{n}\fi}

\newcommand{\Lean}{\textsc{Lean}}
\ExplSyntaxOn
\NewDocumentCommand{\leanid}{m}{
  \ifmmode
    \textsf{\detokenize{#1}}
  \else
    \group_begin:
    \urlstyle{sf}
    \tl_set:Nn \l_tmpa_tl {#1}
    \tl_replace_all:Nnn \l_tmpa_tl {\_} {_}
    \exp_args:NV \path \l_tmpa_tl
    \group_end:
  \fi
}
\ExplSyntaxOff
\providecommand{\tr}{\operatorname{tr}}
\newcommand{\tnleanIssueBase}{https://github.com/LionSR/TNLean/issues/}
\newcommand{\tnleanPrBase}{https://github.com/LionSR/TNLean/pull/}
\newcommand{\tnleanDocBase}{https://sirui-lu.com/TNLean/docs/}
\newcommand{\ghissue}[1]{\href{\tnleanIssueBase#1}{GitHub issue \##1}}
\newcommand{\ghpr}[1]{\href{\tnleanPrBase#1}{PR \##1}}
\newcommand{\doclink}[1]{\href{\tnleanDocBase#1.html}{\path{#1}}}

% Dirac notation and the identity operator, as in blueprint/src/macros/common.tex.
% \providecommand keeps note-local definitions authoritative.
\usepackage{dsfont}
\providecommand{\ket}[1]{|#1\rangle}
\providecommand{\bra}[1]{\langle #1|}
\providecommand{\braket}[2]{\langle #1 | #2 \rangle}
\providecommand{\ketbra}[2]{|#1\rangle\!\langle #2|}
\providecommand{\Id}{\mathds{1}}
\providecommand{\one}{\mathds{1}}
\providecommand{\C}{\mathbb{C}}
\providecommand{\MN}[1]{M_{#1}(\C)}
\providecommand{\MD}{M_D(\C)}

% Verdict marker: \gapnote{<kind>}{<status>}. Kind is the policy
% classification (clarification, local-correction, scope-restriction,
% unfaithful, false-source, open-gap); status is open, wip (elimination
% underway), resolved, or historical. Machine-read by the site index;
% prints a verdict line.
\newcommand{\gapnote}[2]{\par\noindent\textbf{Verdict:} #1 (#2).\par}


% Fallbacks keep this note compilable when a different command.tex is found first.
\providecommand{\leanid}[1]{\nolinkurl{#1}}
\providecommand{\gapnote}[2]{\par\noindent\textbf{Verdict:} #1 (#2).\par}

\title{A Basis of a Degenerate Ground Space in Hastings' Local Solution
  of Gapped Hamiltonians}
\author{}
\date{2026-10-24}

\begin{document}
\maketitle
\gapnote{scope-restriction}{open}

Hastings, \emph{Solving gapped Hamiltonians locally}~\cite{Hastings2006Solving},
\href{https://arxiv.org/abs/cond-mat/0508554v3}{arXiv:cond-mat/0508554v3},
source \path{proj3new.tex}, lines 1678--1705, treats a degenerate ground space
by choosing a basis in which the compressions of the filtered terms are nearly
diagonal. The filtered terms $M_i$, the projection $P$ onto the ground space
and the low-energy commutators $[PM_iP,PM_jP]$ are those of the note
\path{hastings05_filtered_term_approximation.tex} on the filtered terms.

\section{The source passage and what is proved}

For a ground space of dimension $n=2$ the source chooses a basis that
diagonalizes $A=PM_jP$, with eigenvalues $r\ne s$ and
$\delta\lambda_j=|r-s|$ (lines~1678--1688). For an orthonormal eigenbasis
$v,w$ and any $B$, $\langle v,[A,B]w\rangle=(r-s)\langle v,Bw\rangle$, so
both off-diagonal elements of $B$ are at most $\lVert[A,B]\rVert/|r-s|$. The
off-diagonal part
$\langle v,Bw\rangle\lvert v\rangle\langle w\rvert
+\langle w,Bv\rangle\lvert w\rangle\langle v\rvert$ maps the two
coordinates of a vector to orthogonal directions, so its norm is at most the
larger element, and
\begin{align}
  o(i)\le\frac{\lVert[PM_iP,PM_jP]\rVert}{\delta\lambda_j}
  \label{eq:hastings05_two_dim_ground_basis_oi}
\end{align}
(lines~1688--1690). Since $PM_iv$ lies in the span of $v,w$ and
$\langle v,PM_iv\rangle=\langle v,M_iv\rangle$,
\begin{align}
  \lVert M_iv-\langle v,M_iv\rangle v\rVert
  \le\lVert(1-P)M_iv\rVert+o(i),
  \label{eq:hastings05_two_dim_ground_basis_low}
\end{align}
which is eq.~\texttt{low} up to $o(i)$ (lines~476--491); eq.~\texttt{facld}
bounds the first term.

\section{The restriction}

Every dimension $n$ is treated, by the source's iteration over further
compressions with any number of levels (the last two sections), ending when
every pair of basis vectors is separated, as the source's iteration ends. The basis is taken as given, as the
source chooses it, and the $j$ maximizing $\delta\lambda_j$ is not chosen,
since the bounds hold for every $j$. One restriction remains: the
formalization leaves eq.~\texttt{oibu} as
$\sum_io(i)\le\sum_i\lVert[PM_iP,PM_jP]\rVert/\delta\lambda_j$: the
source sums the decaying commutator bound over $i$ only to order of magnitude,
and the resulting constant depends on the growth of the graph.

\section{A compression with simple spectrum}

For $n>2$ the source finds the $M_j$ with the largest spread of eigenvalues,
bounds the off-diagonal elements between its eigenvectors, and passes to
further $M_k$ for eigenvectors sharing an eigenvalue of $M_j$
(lines~1695--1705). Because the compressions commute only approximately, this
iteration is not a proof as stated. When one compression $A=PM_jP$ has $n$
eigenvalues on the ground space that are pairwise at least $\delta>0$ apart,
its orthonormal eigenbasis $v_0,\dots,v_{n-1}$ needs no iteration. With indices
modulo $n$, the off-diagonal part $o_B$ of any $B$ is the sum over $k\ne0$ of
the shifts
$S_k=\sum_a\langle v_a,Bv_{a+k}\rangle\lvert v_a\rangle\langle v_{a+k}\rvert$.
By Bessel's inequality each shift has norm at most its largest element, and
each element is at most $\lVert[A,B]\rVert/\delta$, so
\begin{align}
  o(i)\le(n-1)\frac{\lVert[PM_iP,PM_jP]\rVert}{\delta},
  \label{eq:hastings05_two_dim_ground_basis_simple}
\end{align}
which for $n=2$ is~\eqref{eq:hastings05_two_dim_ground_basis_oi}. The bound
\eqref{eq:hastings05_two_dim_ground_basis_low} holds for every basis vector
$v_a$.\footnote{\leanid{TNLean.Lattice.norm_offDiagonalPartFamily_le_norm_commutator_div},
    \leanid{TNLean.Lattice.sum_norm_offDiagonalPartFamily_le} and
    \leanid{TNLean.Lattice.norm_apply_sub_inner_smul_le_add_offDiagonalPartFamily}
    in \doclink{TNLean/Lattice/LowEnergyBasisSimpleSpectrum}.}
This is the first level of the iteration of the next section, with $\delta$ a
given lower bound on the spacing.

\section{Repeated eigenvalues: the iteration}

Record the source's iteration by real labels $r_\ell(a)$ of the basis vectors at
levels $\ell=0,1,\dots$: the class of $a$ at level $\ell$ is the set of $c$
with $r_m(c)=r_m(a)$ for every $m<\ell$, and the compression $A_\ell=PM_{j_\ell}P$
of level $\ell$ acts diagonally on each class,
$\langle v_c,A_\ell v_d\rangle=r_\ell(c)\delta_{cd}$. Expanding both products in
the basis gives, for $a$, $b$ in one class $C$ at level $\ell$,
\begin{align}
  (r_\ell(a)-r_\ell(b))\langle v_a,Bv_b\rangle
  =\langle v_a,[A_\ell,B]v_b\rangle
  -\sum_{c\notin C}\bigl(\langle v_a,A_\ell v_c\rangle\langle v_c,Bv_b\rangle
    -\langle v_a,Bv_c\rangle\langle v_c,A_\ell v_b\rangle\bigr).
  \label{eq:hastings05_two_dim_ground_basis_iteration}
\end{align}
The indices $c\notin C$ are separated from $a$ and $b$ at an earlier level, so
the sum is a product of matrix elements already bounded there. Dividing by the
smallest difference $\delta_\ell>0$ of distinct labels of level $\ell$ within a
class, which is positive in finite dimension and is the source's
$\delta\lambda$ at that level, bounds every off-diagonal element at the first
level separating its two vectors. With two levels, $M_j$ and then $M_k$, every
off-diagonal element of $PM_iP$ is at most
\begin{align}
  \max\Bigl(\frac{\lVert[A_0,B]\rVert}{\delta_0},\
  \frac{1}{\delta_1}\Bigl(\lVert[A_1,B]\rVert
    +2n\frac{\lVert[A_0,A_1]\rVert\,\lVert[A_0,B]\rVert}{\delta_0^2}\Bigr)\Bigr),
  \label{eq:hastings05_two_dim_ground_basis_two_levels}
\end{align}
and $o(i)$ is at most $n-1$ times this. For any number of levels, let every
commutator $[A_\ell,A_m]$ be at most $K$ and every $[A_\ell,B]$ at most $K_B$.
By induction on the level, the elements of the $A_m$ separated at level $\ell$
are at most $\alpha_\ell$, and those of $B$ at most $\beta_\ell$, where
$\alpha_0=K/\delta_0$, $\beta_0=K_B/\delta_0$ and
\begin{align}
  \alpha_{\ell+1}=\max\Bigl(\alpha_\ell,\frac{K+2n\alpha_\ell^2}{\delta_{\ell+1}}\Bigr),
  \qquad
  \beta_{\ell+1}=\max\Bigl(\beta_\ell,\frac{K_B+2n\alpha_\ell\beta_\ell}{\delta_{\ell+1}}\Bigr).
  \label{eq:hastings05_two_dim_ground_basis_levels}
\end{align}
The source's iteration ends once every off-diagonal element is bounded
(lines~1707--1708), that is, once every pair of basis vectors is separated at
some level; after $L$ such levels, $o(i)\le(n-1)\beta_{L-1}$.\footnote{\leanid{TNLean.Lattice.abs_sub_mul_norm_inner_le},
    \leanid{TNLean.Lattice.norm_inner_le_of_mem_labelClass},
    \leanid{TNLean.Lattice.norm_inner_le_levelBoundB},
    \leanid{TNLean.Lattice.norm_offDiagonalPartFamily_le_levelBoundB},
    \leanid{TNLean.Lattice.norm_inner_le_of_two_levels} and
    \leanid{TNLean.Lattice.norm_offDiagonalPartFamily_le_of_two_levels} in
    \doclink{TNLean/Lattice/LowEnergyBasisIteration}.}
The bound $\beta_\ell$ is linear in $K_B$, and with two levels every term decays
in $d(i,j)$ or $d(i,k)$, so the sum over $i$ behind eq.~\texttt{oibu} stays
bounded.

This is a local correction of the sketch. The source states that the further
compressions bound the remaining off-diagonal elements, but the commutator of
the compression of $M_k$ to a class with that of $B$ differs from $[M_k,B]$ by
the cross terms of~\eqref{eq:hastings05_two_dim_ground_basis_iteration}, which
the sketch omits and which are kept here.

Elimination plan for the remaining restriction: sum the commutator bound over
$i$ with the growth constants of the graph, giving eq.~\texttt{oibu}
explicitly.

\section{Formalization}

\leanid{TNLean.Lattice.norm_offDiagonalPart_le_norm_commutator_div} and
\leanid{TNLean.Lattice.sum_norm_offDiagonalPart_le}
prove~\eqref{eq:hastings05_two_dim_ground_basis_oi} and its sum, and
\leanid{TNLean.Lattice.norm_apply_sub_inner_smul_le_add_offDiagonalPart}
proves~\eqref{eq:hastings05_two_dim_ground_basis_low}, all in
\doclink{TNLean/Lattice/LowEnergyBasis}.

\bibliographystyle{plain}
\bibliography{references}
\end{document}
